\subsection{The subgroup generated by the elliptic period locus}

\begin{result}{A.4.9}{The elliptic period locus generates the full motivic torsor}{Known; translation proved}
For every elliptic curve over $\Qb$, the smallest algebraic subgroup containing a translated period locus is the full motivic group, in the category specified below.
\end{result}
\begin{proof}
We fix the displayed inclusion \(F=\overline{\mathbf Q}\subset\mathbf C\),
an elliptic curve \(E/F\), and the neutral Tannakian category \(\langle
h^1(E)\rangle\) inside André's motivated motives over \(F\), with rational
coefficients. Its Betti fibre functor is \(B\), and its de Rham fibre functor
is \(D\). Put \(G=\operatorname{Aut}^{\otimes}(B)\), a group over
\(\mathbf Q\), and
\[
P=\operatorname{Isom}^{\otimes}(B\otimes_{\mathbf Q}F,D).
\]
This is a right \(G_F\)-torsor: the action is composition on the source.
Write \(\alpha\in P(\mathbf C)\) for comparison in this direction, and let
\(Z\subseteq P\) be its reduced \(F\)-Zariski closure. Choose
\(z_0\in Z(F)\). Such a point exists because \(Z\) is nonempty and of
finite type over the algebraically closed field \(F\). Define \(H\) to be
the smallest closed algebraic subgroup of \(G_F\) containing the image of
\[
 Z\longrightarrow G_F,\qquad z\longmapsto z_0^{-1}\circ z.
\]
Then
\[
H=G_F.
\]
For a curve without CM this is \(\operatorname{GL}_{2,F}\). For a CM curve
with endomorphism field \(K=\operatorname{End}(E)\otimes\mathbf Q\), it is
\((\operatorname{Res}_{K/\mathbf Q}\mathbf G_m)_F\).

\textbf{Proof of the translation and derivation.}

We first verify exactly what subgroup generation means. A closed nonempty
subscheme \(S\subseteq P\) that is stable under the ternary operation
\((x,y,z)\mapsto x\circ y^{-1}\circ z\) and contains \(z_0\) is of the
form \(z_0H_S\), with \(H_S=z_0^{-1}S\) a closed subgroup of \(G_F\).
Indeed, using \(z_0\) in the middle slot proves closure of \(H_S\) under
multiplication, and using it in the first and last slots proves closure
under inversion. Conversely a translate of a subgroup has this ternary
closure. Consequently \(z_0H\) is exactly the smallest algebraic subtorsor
of \(P\) containing \(Z\).

This construction does not depend on the chosen point of \(Z\). For
\(z_1\in Z(F)\), put \(a=z_0^{-1}z_1\in H(F)\), and let
\(S=z_0^{-1}Z\). The set \(a^{-1}S=z_1^{-1}Z\) lies in \(H\), contains
\(a^{-1}\) because \(1\in S\), and therefore its generated subgroup
contains \(a\) and every element \(a(a^{-1}s)=s\) of \(S\). Its generated
subgroup is thus exactly \(H\).

Bost--Charles, Definition 2.3, calls the smallest subtorsor containing the
period locus \(\Omega_E\); their Definition 2.9 defines the motivic torsor
\(\Omega_E^{\mathrm{mot}}\). Their comparison is written
\(c:D_{\mathbf C}\to B_{\mathbf C}\), so the map matching their torsor
with ours is inversion \(c\mapsto c^{-1}\). This is an algebraic
isomorphism of the two isomorphism schemes, carries their comparison point
to \(\alpha\), and carries their ternary operation to the same ternary
operation with the first and third arguments reversed. It therefore
preserves the property of being the generated subtorsor.

The precise published-preprint input is Kreutz--Shen--Vial,
\emph{De Rham--Betti classes with coefficients}, arXiv:2206.08618v3 (21 January
2026), Proposition 3.1 and Theorem 5.13(ii). Proposition 3.1 identifies
\(\Omega_E\) with the Tannakian de Rham--Betti torsor after extending Betti
coefficients to \(F\). Theorem 5.13(ii), applied with one elliptic factor
and coefficient field \(L=F\), identifies that torsor with the motivic
torsor. Its restriction on the number of CM fields is satisfied: a single
elliptic curve contributes zero or one such field.

Their statement is for \(h(E)\). Here the categories generated by \(h(E)\)
and by \(h^1(E)\) are equal, since
\[
h(E)=\mathbf1\oplus h^1(E)\oplus\bigwedge^2h^1(E),
\qquad \bigwedge^2h^1(E)=\mathbf Q(-1).
\]
The Künneth projectors and the identification of the last summand are
algebraic, using the rational point which is the origin of \(E\). An
isomorphism on \(h^1(E)\) extends to \(h(E)\) by the identity on the unit
and the exterior square on \(h^2(E)\), while restriction to the direct
summand is its inverse. These are algebraic maps on comparison torsors.
Thus the cited equality applies to the exact torsor \(P\) above and gives
\(z_0H=P=z_0G_F\), hence \(H=G_F\).

The group description follows from KSV Proposition 5.14 for a single
elliptic curve: its Hodge group is \(\operatorname{SL}_2\) in the non-CM
case and the norm-one torus of \(K\) in the CM case. Including the weight
homotheties gives respectively \(\operatorname{GL}_2\) and
\(\operatorname{Res}_{K/\mathbf Q}\mathbf G_m\). The same proposition
states algebraicity of Hodge classes on all powers, so these groups also
give the motivic group with algebraic correspondences. This completes the
derivation in the manuscript's notation.

\textbf{Relation to the full period conjecture, with an exact morphism.}

The map in question is the inclusion of the period locus into its
generated torsor,
\[
Z\hookrightarrow z_0\langle z_0^{-1}Z\rangle=P.
\]
A.4.9 proves equality of the generated torsor with \(P\), and it does not
identify this displayed inclusion with an isomorphism. Its replacement by
an isomorphism is exactly Bost--Charles Conjecture 2.12. A general-torsor
example demonstrating the difference is A.5.2 below. No counterexample to
the elliptic period conjecture is claimed.

\textbf{Role of the analytic subgroup theorem.}

KSV Proposition 5.1 states the analytic subgroup theorem in the category
generated additively by first cohomology of abelian varieties, with exact
reference to Huber--Wüstholz Theorem 9.7. KSV Proposition 5.8 proves its
version with algebraic coefficient field \(L\); Theorem 5.10 derives full
faithfulness and stability under subobjects for that additive category;
Theorem 5.12 derives connectedness, reductivity, the endomorphism
centralizer, and the nontrivial determinant of the de Rham--Betti group.
For a non-CM elliptic curve these properties force the group to be
\(\operatorname{GL}_2\), as used in Theorem 5.13. In the CM case the proof
of 5.13(ii) explicitly uses Chudnovsky. This chain is the verified way that
the analytic subgroup theorem enters the present result. We do not replace
the coefficient-field argument by an unproved identification of a
\(\mathbf Q\)-Tannakian group with a subgroup defined only over
\(\overline{\mathbf Q}\).

The primary theorem locators used in this argument are \cite[Proposition 3.1; Theorem 5.13(ii); Proposition 5.14]{GenKSV2026}.
\end{proof}

\subsection{The automorphism image and period genericity}

\begin{result}{A.5.1}{Exact orbit-closure formula for the cocycle of 18:59--61}{Proved}
The actual automorphism-cocycle image is Zariski dense exactly when the comparison point is generic over the same base field.
\end{result}
\begin{proof}
Let \(F\) be any characteristic-zero subfield of \(\mathbf C\), let
\(G/F\) be an affine algebraic group of finite type, let \(P/F\) be a
right \(G\)-torsor, and let \(p\in P(\mathbf C)\). Write
\(\Gamma=\operatorname{Aut}(\mathbf C/F)\), and define
\[
\varphi_p:\Gamma\longrightarrow G(\mathbf C),\qquad
\varphi_p(\sigma)=(\sigma p)^{-1}p.
\]
The notation \(x^{-1}y\) is the unique \(a\in G\) with \(y=x\cdot a\);
it is a morphism \(P\times_F P\to G\). The group \(\Gamma\) may be
given its pointwise-stabilizer topology (stabilizers of finite subsets of
\(\mathbf C\) form an identity-neighborhood basis). No topological
continuity assertion is used. All closures here are Zariski closures over
\(\mathbf C\).

Let \(Z_F(p)\) denote the reduced \(F\)-Zariski locus of \(p\) in
\(P\). Then
\[
\overline{\operatorname{im}\varphi_p}^{\mathrm{Zar}}
=D_p\bigl((Z_F(p))_{\mathbf C}\bigr),\qquad
D_p:P_{\mathbf C}\overset\sim\longrightarrow G_{\mathbf C},
\quad x\longmapsto x^{-1}p.
\tag{A5.1.1}
\]
The inverse morphism is \(a\mapsto p\cdot a^{-1}\). In particular,
\[
\overline{\operatorname{im}\varphi_p}^{\mathrm{Zar}}=G_{\mathbf C}
\quad\Longleftrightarrow\quad Z_F(p)=P,
\tag{A5.1.2}
\]
and its dimension is \(\operatorname{trdeg}_F F(p)\).

\textbf{Proof.}

We first show \(\mathbf C^{\Gamma}=F\). An element algebraic over
\(F\) but outside \(F\) has a distinct conjugate in characteristic zero,
and the corresponding \(F\)-embedding of its finite extension extends to
an automorphism of \(\mathbf C\): first extend to algebraic closures, then
to transcendence bases and algebraic closures again. A transcendental
element \(t\) can be moved by the \(F\)-automorphism \(t\mapsto t+1\) of
\(F(t)\), extended to \(\mathbf C\) by the same argument. Thus precisely
the elements of \(F\) are fixed by every member of \(\Gamma\).

Because a torsor under an affine group is affine, embed \(P\) into an
affine space over \(F\). Let \(S=\Gamma p\), and let \(I\) be the ideal
of its \(\mathbf C\)-Zariski closure in the polynomial ring of that affine
space. The semilinear action of \(\Gamma\) preserves \(I\), since it
preserves \(S\). For each total-degree bound \(d\), the finite-dimensional
space \(I_{\le d}\) is stable under coefficientwise action of \(\Gamma\).
Write its unique reduced row-echelon basis in the ordered monomial basis
of polynomials of degree at most \(d\). Applying any \(\sigma\in\Gamma\)
to this basis gives another reduced row-echelon basis for the same
subspace. Uniqueness fixes every coefficient; the fixed-field statement
therefore puts them all in \(F\). Hence \(I\) is generated over
\(\mathbf C\) by its intersection with \(F[x_1,\ldots,x_n]\).

For a polynomial with coefficients in \(F\), vanishing on the whole orbit
\(S\) is equivalent to vanishing at \(p\), because
\(f(\sigma p)=\sigma(f(p))\). Thus that intersection is exactly the ideal
of \(Z_F(p)\). We have proved
\(\overline{S}^{\mathrm{Zar}}=(Z_F(p))_{\mathbf C}\). Applying the
algebraic isomorphism \(D_p\), which takes \(\sigma p\) to
\(\varphi_p(\sigma)\), proves (A5.1.1). Faithful base change proves the
equivalence (A5.1.2). Finally, the coordinate algebra of the locus embeds
in \(\mathbf C\) by evaluation at \(p\), has fraction field \(F(p)\),
and is a finite-type integral \(F\)-algebra. Its Krull dimension is the
displayed transcendence degree. Base change and the isomorphism \(D_p\)
preserve the dimension.

\textbf{Application to the page and to the named conjecture.}

Scan 18:59 says: ``ce point \ldots{} est générique'' and then ``a une image Zariski
dense''. These two formulations are exactly equivalent by A.5.1, once
\(P\) is identified with the isomorphism torsor of the two fibre functors
and \(\varphi\) is the corrected comparison cocycle. In particular this
is a formal equivalence and does not prove the claimed genericity.

For a motive over \(\overline{\mathbf Q}\), choose \(F=\overline{
\mathbf Q}\), \(P\) the motivic comparison torsor, and \(p\) the period
comparison. Bost--Charles Conjecture 2.12 is the statement
\(Z_F(p)=P\). Inversion of the comparison direction identifies their
torsor with the one in the manuscript, as in A.4.9, and A5.1.2 then proves
the requested equivalence with density of the \textbf{actual image} of
\(\varphi\).

The page's ambient field \(k\simeq\mathbf C\), its finite-type
definition field \(k_0\), and its remark about minimum transcendence degree
also allow transcendental \(k_0\). A.5.1 remains valid for that data. The
statement ``the point is generic'' over such \(k_0\) is an additional
assertion, whereas the explicit Bost--Charles conjecture being compared
has base \(\overline{\mathbf Q}\). We have proved the exact equivalence
over a fixed base field; no transfer of genericity to a new transcendental
base field is used or implied. For an actual comparison point, adjoining
one of its transcendental coordinates to the base can lower its relative
transcendence degree. Therefore an application must preserve the field
named in the statement.

There is also no need to trivialize \(P\) over a finite extension of
\(F\) for this proof. This matters for nonconnected \(G\): a finite
extension can split components, and a single point need not be dense in
their union. The intrinsic difference map \(P\times P\to G\) proves
the formula without that change of base.

The primary theorem locators used in this argument are \cite[Definitions 2.3 and 2.9; Conjecture 2.12]{GenBostCharles}.
\end{proof}

\begin{result}{A.5.2}{Actual image and generated subgroup have different closures}{Proved counterexample}
Passing from the actual cocycle image to its generated subgroup can change its Zariski closure.
\end{result}
\begin{proof}
Take \(F=\mathbf Q\), \(G=P=\mathbf G_m\),
\(g=1+\sqrt2\in\mathbf C^\times\), and
\(\Gamma=\operatorname{Aut}(\mathbf C/\mathbf Q)\). The two conjugates
of \(g\) are \(1+\sqrt2\) and \(1-\sqrt2\), so
\[
\operatorname{im}(\sigma\mapsto\sigma(g)^{-1}g)
=\{1,-3-2\sqrt2\}.
\]
This is a closed zero-dimensional subset of \(\mathbf G_{m,\mathbf C}\).
The subgroup it generates contains all integer powers of
\(r=-3-2\sqrt2\). Since \(|r|=3+2\sqrt2>1\), these are pairwise
distinct. A nonzero Laurent polynomial has finitely many roots in
\(\mathbf C^\times\), after multiplication by a monomial reduces it to
an ordinary nonzero polynomial. The subgroup is therefore Zariski dense.
The exact set inclusion
\[
\operatorname{im}\varphi\ \subset\ \langle\operatorname{im}\varphi\rangle
\]
thus has zero-dimensional closure on its left and one-dimensional closure
on its right. This verifies P26 {[}II{]} and proves why the word ``image'' must
be retained in the equivalence of A.5.1.

The calculation uses exact SymPy arithmetic; its script is
\path{code/check_generation.py}, with output in
\path{code/check_generation.out}.
the infinitude and Zariski-density arguments are the proofs above. The
counterexample does not have \(F=\overline{\mathbf Q}\), and no motive
with these as its prescribed comparison coordinates and motivic group is
constructed. Consequently it does not refute any motivic period
conjecture; it refutes only the general formal substitution of generated
subgroup density for actual-image density.
\end{proof}

\begin{result}{A.5.3}{Cocycle conventions and the semidirect-product homomorphism}{Proved}
Inversion converts the displayed right cocycle to a left cocycle; its graph is a homomorphism into the explicitly defined semidirect product.
\end{result}
\begin{proof}
In the data of A.5.1, composition in \(\Gamma\) is
\((\tau\sigma)(x)=\tau(\sigma(x))\). Direct composition of the typed
isomorphisms gives
\[
\begin{aligned}
{}^\tau\varphi_p(\sigma)\,\varphi_p(\tau)
&=((\tau\sigma)p)^{-1}(\tau p)\,(\tau p)^{-1}p\\
&=((\tau\sigma)p)^{-1}p
=\varphi_p(\tau\sigma).
\end{aligned}
\]
Thus \(\varphi_p\) obeys the page's right-cocycle convention. Pointwise
inversion \(\psi_p(\sigma)=\varphi_p(\sigma)^{-1}\) gives
\[
\psi_p(\tau\sigma)=\psi_p(\tau)\,{}^\tau\psi_p(\sigma).
\]
For the group law
\((a,\tau)(b,\sigma)=(a\,{}^\tau b,\tau\sigma)\) on
\(G(\mathbf C)\rtimes\Gamma\), the map
\(\sigma\mapsto(\psi_p(\sigma),\sigma)\) is therefore a group
homomorphism. Every source and target here is a group of points, rather
than an assertion about continuity in an adelic topology.

After an \(F\)-rational trivialization \(P\simeq G\), with period point
represented by \(g\), the formula becomes
\(\varphi_p(\sigma)=({}^\sigma g)^{-1}g\), exactly the live correction
on 18:61. For \(G=\mathbf G_m\) and a transcendental \(t\), choose
commuting automorphisms \(\sigma(t)=t+1\), \(\tau(t)=t+2\), for example
\(\tau=\sigma^2\) after extending \(\sigma\) to \(\mathbf C\). Then
\[
\varphi_p(\tau\sigma)=\frac{t}{t+3},\quad
\varphi_p(\tau)\varphi_p(\sigma)=\frac{t^2}{(t+2)(t+1)}.
\]
Their difference (ordinary product minus composite) is
\(-2t/((t+1)(t+2)(t+3))\ne0\). The semilinear product, however, has
zero residual. The exact symbolic computation is in the same script and
output. This supplies the specific group morphism to which the page's
``hom total'' can be compared, without asserting that the archimedean
cocycle alone is a homomorphism to \(G(\mathbf C)\).
\end{proof}

\begin{result}{A.5.4}{A genuine motivic counterexample over a transcendental field of definition}{Proved counterexample}
There are genuine elliptic period comparisons over a transcendental field for which the actual image is not dense but its generated subgroup is dense. An augmented degree-one example has minimum field transcendence degree one and still has strictly different closures.
\end{result}
\begin{proof}
This supplements A.5.2 with
an actual comparison point of an abelian motive. It also refutes a universal
reading of the 18:59 genericity assertion over arbitrary finitely generated
fields of definition, even after imposing the written qualification that
\(\pi\) be transcendental over the base. It does not concern the arithmetic
base \(\overline{\mathbf Q}\) of Bost--Charles Conjecture 2.12.

Take \(E_0:y^2=x^3-x\). Its CM field is \(K=\mathbf Q(i)\). Enlarge
\(K\) by a finite extension \(K_1\) so that the CM period torsor has an
algebraic base point and the chosen split eigenbases are defined over
\(K_1\). This is possible because it is a nonempty finite-type torsor
over a number field; a closed point has finite residue field. In these
bases its comparison point has coordinates \((a,b)\) in
\((\mathbf G_m)^2\). Choose the first eigenspace to be the one arising
from the holomorphic differential \(dx/y\). Then \(a\) is a nonzero
algebraic multiple of a nonzero holomorphic period, or its inverse if the
comparison is inverted. The CM field calculation of A.4.3, using Chudnovsky \cite[\S2A, Theorem 1, p.343]{EllChudnovsky1978}, gives
\[
\operatorname{trdeg}_{K_1}K_1(a,\pi)=2.
\]
Inverting \(a\) and multiplying it by a nonzero algebraic number preserve
this statement by the explicit inverse field maps.

Put
\[
s=a+a^{-1},\qquad F=K_1(s)\subset\mathbf C.
\]
The inclusion \(F\subset K_1(a)\) has degree two. Indeed, the involution
\(a\mapsto a^{-1}\) of the rational function field \(K_1(a)\) is
nontrivial and fixes \(s\), while \(a\) satisfies
\(X^2-sX+1=0\). Moreover \(\pi\) is transcendental over \(F\), since
it is algebraically independent of \(a\) over \(K_1\).

Already for the actual motive \(h^1(E_0)_F\), this gives a complete
counterexample to replacing density of the image by density of its
generated subgroup. Its split comparison coordinates satisfy
\(ab=\lambda(2\pi i)^{\epsilon}\), with
\(\lambda\in K_1^\times\) and \(\epsilon=1\) or \(-1\) according to
the comparison direction. Thus \(b\) is transcendental over
\(K_1(a)\). Automorphisms fixing \(F(a)\) can send \(b\) to any
element transcendental over \(F(a)\), so the closure of their cocycle
values is \(\{1\}\times\mathbf G_m\). An automorphism sending
\(a\) to \(a^{-1}\) supplies a second coset. Consequently the full
actual-image closure is
\[
\{1,a^2\}\times\mathbf G_m\subsetneq T_{\mathbf C},
\]
whereas its generated subgroup is Zariski dense in
\(T_{\mathbf C}=\mathbf G_m^2\). This CM motive descends to a number
field, so \(F\) is not a minimum-transcendence-degree field for this
first example. The construction below incorporates a varying elliptic
factor to address that additional quantifier; it establishes strictness
of the two closures there without asserting that the larger one is the
entire group.

An explicit elliptic curve over \(F\) with \(j\)-invariant \(s\) is
\[
E_s:\quad y^2=x^3-\frac{3s}{s-1728}x-\frac{2s}{s-1728}.
\]
Substitution in \(j=1728\,4A^3/(4A^3+27B^2)\) gives \(j(E_s)=s\).
Both denominators and the discriminant are nonzero because \(s\) is
transcendental. Thus \(E_s\) has no CM and is not isogenous to \(E_0\):
CM is preserved by isogeny, and a CM elliptic curve has algebraic
\(j\)-invariant. Let
\[
M=h^1(E_0)\oplus h^1(E_s)
\]
in the rational category of abelian motives, and form its tensor
comparison torsor \(P_M/F\). Restriction to the first summand gives a
surjective morphism of motivic groups \(G_M\to T\), where
\(T=\operatorname{Res}_{K/\mathbf Q}\mathbf G_m\). One can see the
surjectivity directly from Tannaka duality: the subcategory generated by
the first summand is full and closed under subobjects. This is
Deligne--Milne Proposition 2.21(a). After extension to \(F\), compose this
map with the first split character \(T_F\to\mathbf G_{m,F}\), obtaining
\(\chi:G_{M,F}\to\mathbf G_{m,F}\).

Apply the same character to the comparison torsor, using the fixed
algebraic base point on the CM factor. It gives a regular invertible
function \(f\) on \(P_M\) with \(f(p_M)=a\); it is nonconstant because
\(\chi\) is surjective. For the cocycle convention of A.5.1,
\[
\chi(\varphi_{p_M}(\sigma))
=\frac{a}{\sigma(a)}\in\{1,a^2\}.
\tag{A.5.4.1}
\]
Both values occur. The nontrivial automorphism of \(K_1(a)/F\) extends
first to its algebraic closure and then, after choosing a transcendence
basis, to an automorphism of \(\mathbf C/F\).

Let \(Y\) be the complex Zariski closure of the actual cocycle image,
and let \(J\) be the closed algebraic subgroup generated by that image.
Then
\[
Y\subseteq\chi^{-1}(\{1,a^2\}).
\]
Choose \(u\) in the image with \(\chi(u)=a^2\). Its square belongs to
\(J\) and has character value \(a^4\), which belongs to neither
\(\{1,a^2\}\) because \(a\) is transcendental. Therefore
\[
Y\subsetneq J.
\tag{A.5.4.2}
\]
Indeed \(\chi(J)\) contains all powers of \(a^2\), so its Zariski
closure is \(\mathbf G_m\). This proves a strict change in the density
answer under the actual-image-to-generated-subgroup map using a genuine
motivic comparison point. It does not assert that \(J=G_M\), which is
unneeded for the counterexample.

The comparison point itself is not generic over \(F\): it satisfies the
nonzero equation
\[
f^2-sf+1=0
\]
on \(P_M\). The function is nonzero because \(f\) is surjective onto a
\(\mathbf G_m\)-torsor after trivialization, whereas a nonzero quadratic
polynomial has only finitely many roots.

The field \(F\) has the minimum possible transcendence degree, namely
one, for the abelian variety \(E_0\times E_s\) up to isogeny, and hence
for its degree-one motive in the category of effective degree-one
motives. Here is the precise descent argument. Deligne--Milne
Proposition 6.21 identifies that category with abelian varieties up to
isogeny. If this object descended to a number field, after extending to
\(\overline{\mathbf Q}\) its elliptic isogeny factors would be defined
over \(\overline{\mathbf Q}\); Here is the needed base-extension argument. For abelian varieties \(A,B\) over an algebraically closed field \(L\) and an algebraically closed extension \(L'\), every homomorphism over \(L'\) spreads over an integral finite-type \(L\)-scheme. A homomorphism between abelian varieties is rigid: its deformation fixing the origin is constant, as follows by applying the rigidity lemma to the difference of two fibres. The relative Hom scheme is therefore unramified with discrete fibres, and this family of homomorphisms is constant on a connected base. A closed \(L\)-point of the base gives the original homomorphism over \(L\). Thus \(\operatorname{Hom}_L(A,B)\to\operatorname{Hom}_{L'}(A_{L'},B_{L'})\) is bijective; injectivity is faithful base change. Apply this to the rational idempotent cutting out the elliptic factor, clear its denominator, and take the image of the resulting algebraic homomorphism. This produces the elliptic isogeny factor over \(L=\overline{\mathbf Q}\). The factor isogenous to \(E_s\) would therefore be
defined over \(\overline{\mathbf Q}\). The kernel of an isogeny from
an elliptic curve over \(\overline{\mathbf Q}\) is a finite subgroup of
its torsion and is also defined over \(\overline{\mathbf Q}\); its
quotient has algebraic \(j\). This contradicts \(j(E_s)=s\).
The minimum statement is made in this specified degree-one category;
no assertion about a different unspecified notion of descent is needed.

Consequently the word ``minimum'' on 18:53 and the unfinished \(\pi\)
qualification on 18:59 do not, by themselves, repair the universal
genericity assertion for such fields. The exact orbit/genericity
equivalence in A.5.1 remains valid: both of its sides are false in this
example. The implication identifying that assertion with the arithmetic
Bost--Charles conjecture must keep the base field
\(\overline{\mathbf Q}\), or supply an additional relative formulation
which removes the period relations already present in \(F\).

The primary theorem locators used in this argument are \cite[Proposition 2.21; Theorem 6.20 and Proposition 6.21]{GenDeligneMilne}.
\end{proof}

